% Work-MAT, hodina 57; priprava pro ucitele, nikoli samostatny zadavaci test.
% Navazuje na pripravy 41--56 a hodinovy rozpis prvniho rocniku.
\documentclass[11pt,a4paper]{article}
% Spolecny styl pro ucitelske pripravy MAT (Work/MAT)
% Verze 1.1, 2026-10-23
\NeedsTeXFormat{LaTeX2e}
\ProvidesPackage{priprava}[2026/10/23 v1.1 Sablona ucitelskych priprav MAT]

% Zakladni typografie a cestina
\RequirePackage[T1]{fontenc}
\RequirePackage[utf8]{inputenc}
\RequirePackage[czech]{babel}
\RequirePackage{lmodern}
\RequirePackage{microtype}
\RequirePackage{amsmath,amssymb}
\RequirePackage{xcolor}
\RequirePackage{enumitem}
\RequirePackage[most]{tcolorbox}
\RequirePackage{tabularx,booktabs,array}
\RequirePackage{fancyhdr}
\RequirePackage{needspace}
\RequirePackage{geometry}
\RequirePackage[hidelinks]{hyperref}

\geometry{a4paper,top=17mm,bottom=18mm,left=23mm,right=23mm}
\definecolor{notegray}{gray}{0.39}
\setlength{\parindent}{0pt}
\setlength{\parskip}{2.5pt plus 1pt minus 0.5pt}
\setlist[itemize]{topsep=2pt,itemsep=1pt,parsep=0pt,leftmargin=1.45em}
\setlist[enumerate]{topsep=2pt,itemsep=2pt,parsep=0pt,leftmargin=1.65em}
\setcounter{secnumdepth}{0}
\renewcommand{\arraystretch}{1.12}

% Zahlavi: vlevo rocnik | globalni tema, vpravo cislo: nazev hodiny
% Zapati zustava zcela prazdne.
\pagestyle{fancy}
\fancyhf{}
\renewcommand{\headrulewidth}{0.35pt}
\renewcommand{\footrulewidth}{0pt}
\setlength{\headheight}{14pt}
\newcommand{\nastavPripravu}[4]{%
  \fancyhead[L]{\small #1\enspace\textbar\enspace #2}%
  \fancyhead[R]{\small #3: #4}%
}

% Minimalisticke cernobile boxy. Kazda definice zvlast.
\tcbset{commonlesson/.style={enhanced,breakable=false,colback=white,colframe=black,
  colbacktitle=white,coltitle=black,boxrule=0.45pt,
  arc=0.8mm,outer arc=0.8mm,left=3mm,right=3mm,top=1.6mm,bottom=1.6mm,
  toptitle=1mm,bottomtitle=0.7mm,fonttitle=\bfseries,
  attach boxed title to top left={xshift=2.5mm,yshift=-1.15mm},
  boxed title style={colback=white,colframe=white,boxrule=0pt,left=1mm,right=1mm}}}
\newtcolorbox{zadaniBox}[1][]{commonlesson,title={Zadání},#1}
\newtcolorbox{prikladBox}[1][]{commonlesson,title={Příklad},#1}
\newtcolorbox{definiceBox}[1][]{commonlesson,title={Definice},#1}

% Spolecny box pro zadani a jeho reseni, bez mezery mezi castmi.
% Pouziti:
%   \begin{prikladReseniBox}
%     Zadani ulohy.
%     \tcblower
%     \begin{align*} ... \end{align*}
%   \end{prikladReseniBox}
% Spodni cast ma jemny sedy podklad a delici linku; lze v ni sazet
% vice radku matematiky. Obe casti maji jen jeden spolecny ramecek.
\tcbset{solutionlesson/.style={
  commonlesson,
  enhanced jigsaw,
  breakable,
  colbacklower=black!3,
  segmentation style={solid,draw=black!25,line width=0.35pt},
  fontlower=\small\color{notegray},
  before lower={\noindent\textbf{Řešení:}\enspace},
  leftlower=3mm,rightlower=3mm
}}
\newtcolorbox{prikladReseniBox}[1][]{solutionlesson,title={Příklad},#1}
\newtcolorbox{zadaniReseniBox}[1][]{solutionlesson,title={Zadání},#1}

% Pomucky pro autora pripravy: sedive poznamky; \reseni pro starsi pripravy
\newcommand{\poznamka}[1]{\par\begingroup\small\color{notegray}\itshape
  \textnormal{\textbf{Poznámka pro učitele.}}\enspace #1\par\endgroup}
\newcommand{\reseni}[1]{\par\begingroup\small\color{notegray}
  \textbf{Očekávané řešení:}\enspace #1\par\endgroup}
\newcommand{\blok}[1]{\needspace{5\baselineskip}\subsection*{#1}}
\newcommand{\vyrok}[1]{\uv{#1}}

\nastavPripravu{1. ročník}{Algebraické výrazy}{57}{Průběžný test a rozbor typických chyb}
\begin{document}
\poznamka{45 minut: 3 min organizace, přibližně 20 min průběžný test podle samostatného schváleného zadání, 3 min odevzdání, 14 min cílená diagnostika a oprava postupů, 5 min závěr. Nezavádíme další učivo. Testové zadání, bodování a klasifikační pravidla nejsou v této přípravě obsaženy.}
\blok{0.--23. minuta: průběžný test}
\begin{zadaniBox}[title={Organizace pro učitele}]
Použij předem připravené a schválené zadání odpovídající skutečně probranému učivu: operace s mnohočleny, algebraické vzorce, dělení, rozklady a lomené výrazy. Před zahájením sděl povolené pomůcky, požadovaný zápis podmínek a způsob odevzdání. Po vymezeném čase všechny práce vyber.
\end{zadaniBox}
\poznamka{Příprava sama \emph{nenahrazuje} žákovský test ani nevytváří další testové varianty. Zadání a bodování spravuj odděleně; pokud by výklad 52--56 ještě nebyl hotov, uprav rozsah testu dle skutečného průběhu výuky.}
\blok{23.--40. minuta: diagnostika bez předjímání výsledků testu}
Bezprostředně po odevzdání lze na tabuli vybrat jednu či dvě z následujících úloh, podle toho, jaké obtíže se projevovaly \emph{při výuce}. Nejde o vyhlášení výsledků právě napsané písemky.
\begin{zadaniReseniBox}[breakable=false,title={Práce s chybou: znaménka a distributivita}]
Žák napsal $-2x(x-3)=-2x^2-6x$. Najdi chybu a vysvětli opravu.
\tcblower
Správně $-2x(x-3)=-2x^2+6x$. Druhý dílčí součin je $(-2x)(-3)=+6x$.
\end{zadaniReseniBox}
\begin{zadaniReseniBox}[breakable=false,title={Práce s chybou: algebraický vzorec}]
Proč není správně $(x-3)^2=x^2-9$? Zapiš správný rozvoj.
\tcblower
Z distributivity dostaneme $x^2-3x-3x+9=x^2-6x+9$. V chybném zápisu chybí prostřední člen a znaménko konstanty.
\end{zadaniReseniBox}
\begin{zadaniReseniBox}[breakable=false,title={Práce s chybou: zkrácení a původní obor}]
Žák napsal $\frac{x^2-9}{x-3}=x+3$ \uv{pro všechna reálná $x$}. Co je potřeba doplnit?
\tcblower
Rozklad $x^2-9=(x-3)(x+3)$ dovoluje krácení jen pro $x\ne3$. Rovnost výrazů platí na původním definičním oboru, nikoliv pro $x=3$.
\end{zadaniReseniBox}
\blok{40.--45. minuta: závěr a příprava na opakování}
\begin{zadaniBox}[title={Krátké sebehodnocení}]
Žáci zformulují jeden postup, který zvládají bezpečně, a jeden typ příkladů, který potřebují procvičit. Připomeň, že příští hodina se soustředí na ověřování správnosti úprav a rozlišení důkazu od zkoušky dosazením.
\end{zadaniBox}
\poznamka{Samostatný domácí úkol po testu není nutný. Pokud je vhodný, zadej dokončení jedné konkrétní chyby z hodiny se zdůvodněním opravy, nikoliv další náhodnou sadu příkladů.}
\end{document}
